\documentclass{situationist-preprint}

\title{Indifference and the Limits of Care}
\msid{PHE-06-b}
\msfamily{PHE-06}
\msdiscipline{Phenomenology}
\mstype{Consequence}
\mskeywords{care; indifference; apathy; everydayness; Heidegger}

\begin{document}
\maketitle

\begin{abstract}
PHE-06-a argues that care is the being of Dasein: every existential is a moment of care or a mode of one. We derive what this implies about indifference. Two kinds must be distinguished: the undifferentiated mode of average everydayness, which Heidegger calls indifference, and affective indifference, in which nothing matters to someone. Both are modes of care, and the second presupposes mattering in order to present things as not mattering. The limit of care is therefore not reached in indifference but only where disclosure itself is absent, and there no mode of Dasein's being is exercised at all.
\end{abstract}

\section{Two indifferences}

Heidegger uses \enquote{indifference} in a technical sense when he introduces average everydayness: Dasein is proximally and for the most part neither authentic nor inauthentic but exists in an undifferentiated mode, which is nonetheless a mode of its being \citep[\S9]{Heidegger1927}. Ordinary language uses the word for something else: the state in which nothing matters to someone, as in apathy, exhaustion or depression. Both have been taken to mark a limit of care, the first because it seems to lack the self-concern that care involves, the second because it seems to lack concern altogether.

\section{The derivation}

\begin{proposition}\label{p:every}
Everyday indifference is a mode of care.
\end{proposition}

\begin{proof}
By Theorem~4.4 of PHE-06-a every existential is a moment of care or a mode of one, and by Lemma~4.3 there none of the three moments can be had without the others. Undifferentiated everydayness is a way of existing as a whole, in which one is ahead of oneself in the possibilities that the \enquote{they} lays out, already in a world, and alongside the things of concern. It realizes all three moments together and is therefore a mode of care, distinguished from authenticity and inauthenticity by whether those possibilities are taken over as one's own, not by whether care is present.
\end{proof}

\begin{proposition}\label{p:affective}
Affective indifference is a privative mode of care, and it presupposes mattering.
\end{proposition}

\begin{proof}
By Principle~P1 of PHE-06-a, disclosure is mattering: to disclose something is to disclose it as mattering or failing to matter in some way. In affective indifference things are disclosed, as the subject's own descriptions show: the room is seen, the friend's voice heard, the task recognized. They are disclosed as not mattering, flat and without pull. Failing to matter is a determination that only something capable of mattering can have, and only to one for whom things can matter. Hence the state is a mode of being alongside and of being ahead of oneself in which possibilities show up as unavailable, which is the privative mode described in the reply to Objection~4 of PHE-06-a.
\end{proof}

Ratcliffe's analysis of existential feelings in depression supports the second proposition from a different direction: the loss of the sense that anything could matter is itself felt, as a change in the space of possibilities within which things are encountered \citep{Ratcliffe2008}. A loss that is felt is a disclosure.

\begin{corollary}\label{c:limit}
Care has a limit only where disclosure is absent, and there there is no mode of Dasein's being, indifferent or otherwise.
\end{corollary}

In dreamless sleep or coma, nothing is disclosed, so nothing is disclosed as indifferent. Such states are not counterexamples to Theorem~4.4 of PHE-06-a, since the theorem concerns the structure of Dasein's being as exercised in disclosure, not every bodily condition of the one who exists. The limit of care is not a further way of being but the absence of any exercise of the being that care describes.

\begin{objection}
Proposition~\ref{p:affective} makes the thesis unfalsifiable: any state in which nothing matters is redescribed as a mode of care.
\end{objection}

\begin{reply}
The thesis would be falsified by a disclosure in which things were given neither as mattering nor as failing to matter, as when an item is registered without any bearing on one's situation. The proposition claims that affective indifference is not such a case, on the ground that its subjects report the absence of mattering as a feature of what they encounter. That claim is answerable to description.
\end{reply}

\section{Conclusion}

Neither everyday nor affective indifference marks a limit of care. The first is the undifferentiated mode in which care proximally exists; the second is a privative mode that presupposes mattering. Care ends only where disclosure ends, and indifference is a way of disclosure.

\begin{thebibliography}{99}\small
\bibitem[Heidegger(1927)]{Heidegger1927} Heidegger, M. (1927). Sein und Zeit. \emph{Jahrbuch für Philosophie und phänomenologische Forschung} 8.
\bibitem[Ratcliffe(2008)]{Ratcliffe2008} Ratcliffe, M. (2008). \emph{Feelings of Being: Phenomenology, Psychiatry and the Sense of Reality}. Oxford University Press.
\end{thebibliography}

\end{document}
